\documentclass[11pt,a4paper]{article}

% ── Packages ──────────────────────────────────────────────────────
\usepackage[utf8]{inputenc}
\usepackage[T1]{fontenc}
\usepackage{lmodern}
\usepackage[margin=2.5cm]{geometry}
\usepackage{amsmath,amssymb}
\usepackage{booktabs}
\usepackage{graphicx}
\usepackage{xcolor}
\usepackage{hyperref}
\usepackage{enumitem}
\usepackage{titlesec}
\usepackage{fancyhdr}
\usepackage{float}
\usepackage{array}
\usepackage{longtable}
\usepackage{tcolorbox}

% ── Colors ────────────────────────────────────────────────────────
\definecolor{accent}{RGB}{31,119,180}
\definecolor{warnbg}{RGB}{255,243,205}
\definecolor{warnborder}{RGB}{255,193,7}
\definecolor{infobg}{RGB}{209,236,241}
\definecolor{infoborder}{RGB}{23,162,184}
\definecolor{dangerbg}{RGB}{248,215,218}
\definecolor{dangerborder}{RGB}{220,53,69}

% ── Hyperref config ───────────────────────────────────────────────
\hypersetup{
  colorlinks=true,
  linkcolor=accent,
  urlcolor=accent,
  citecolor=accent,
  pdftitle={Portfolio Multifractality Impact Analysis},
  pdfauthor={StockParticle Research}
}

% ── Header / Footer ──────────────────────────────────────────────
\pagestyle{fancy}
\fancyhf{}
\fancyhead[L]{\small\textit{Portfolio Multifractality Impact Analysis}}
\fancyhead[R]{\small\textit{StockParticle Research}}
\fancyfoot[C]{\thepage}
\renewcommand{\headrulewidth}{0.4pt}

% ── Custom boxes ─────────────────────────────────────────────────
\newtcolorbox{importantbox}{
  colback=warnbg, colframe=warnborder,
  fonttitle=\bfseries, title=Important,
  boxrule=0.8pt, arc=2pt
}
\newtcolorbox{notebox}{
  colback=infobg, colframe=infoborder,
  fonttitle=\bfseries, title=Note,
  boxrule=0.8pt, arc=2pt
}
\newtcolorbox{warningbox}{
  colback=dangerbg, colframe=dangerborder,
  fonttitle=\bfseries, title=Warning,
  boxrule=0.8pt, arc=2pt
}
\newtcolorbox{tipbox}{
  colback=infobg, colframe=accent,
  fonttitle=\bfseries, title=Tip,
  boxrule=0.8pt, arc=2pt
}

% ── Title section styling ────────────────────────────────────────
\titleformat{\section}{\Large\bfseries\color{accent}}{\thesection}{1em}{}
\titleformat{\subsection}{\large\bfseries}{\thesubsection}{1em}{}
\titleformat{\subsubsection}{\normalsize\bfseries}{\thesubsubsection}{1em}{}

% ══════════════════════════════════════════════════════════════════
\begin{document}

% ── Title Page ───────────────────────────────────────────────────
\begin{titlepage}
\centering
\vspace*{3cm}
{\Huge\bfseries\color{accent} Portfolio Multifractality\\[0.3em] Impact Analysis}\\[1.5cm]
{\Large Findings Report}\\[2cm]
{\large StockParticle Research}\\[1cm]
{\large\today}\\[3cm]

\begin{tabular}{ll}
\toprule
\textbf{Parameter} & \textbf{Value} \\
\midrule
Portfolio & AAPL $\cdot$ JPM $\cdot$ XOM $\cdot$ JNJ $\cdot$ GC=F $\cdot$ BTC-USD $\cdot$ EWJ \\
Weighting & Equal-weight (1/7 each) \\
Data window & $\sim$11.5 years (2014-09-18 $\to$ 2026-04-10) \\
Observations & 2{,}904 aligned daily log-returns \\
$q$ range & $[-5,\; 5]$ step 0.5 \\
Scales & 10--200 days (geometric spacing) \\
Polynomial order & 1 (linear detrending) \\
\bottomrule
\end{tabular}
\end{titlepage}

\tableofcontents
\newpage

% ══════════════════════════════════════════════════════════════════
\section{Baseline Multifractal Characterisation}
% ══════════════════════════════════════════════════════════════════

\subsection{Per-Asset MF-DFA Results}

Table~\ref{tab:baseline} summarises the MF-DFA output for each portfolio component. The generalised Hurst exponent $h(q)$ was estimated across scales $s \in [10, 200]$ for $q \in [-5, 5]$. The spectral width $\Delta h = h(q_{\min}) - h(q_{\max})$ measures the degree of multifractality.

\begin{table}[H]
\centering
\caption{Per-asset MF-DFA baseline results}
\label{tab:baseline}
\begin{tabular}{llccl}
\toprule
\textbf{Asset} & \textbf{Sector} & $\boldsymbol{h(2)}$ & $\boldsymbol{\Delta h}$ & \textbf{Interpretation} \\
\midrule
AAPL    & Tech        & 0.506 & 0.232 & Near-random; moderate multifractality \\
JPM     & Financials  & 0.484 & 0.208 & Slight anti-persistence; moderate \\
XOM     & Energy      & 0.497 & \textbf{0.296} & Near-random; \textbf{high} multifractality \\
JNJ     & Healthcare  & 0.449 & \textbf{0.308} & Anti-persistent; \textbf{highest} $\Delta h$ \\
GC=F    & Commodity   & 0.433 & 0.241 & Anti-persistent; moderate \\
BTC-USD & Crypto      & \textbf{0.566} & \textbf{0.300} & \textbf{Persistent} trends; high \\
EWJ     & Intl ETF    & 0.430 & 0.185 & Anti-persistent; \textbf{lowest} $\Delta h$ \\
\bottomrule
\end{tabular}
\end{table}

\begin{importantbox}
If returns were monofractal (random walk), we would expect $h(2) \approx 0.5$ and $\Delta h \approx 0$ for all assets. Every asset shows significant spectral width, confirming that \textbf{the standard $\sqrt{T}$ VaR scaling assumption is fundamentally invalid} for this universe.
\end{importantbox}

\subsection{Portfolio-Level MF-DFA}

\begin{table}[H]
\centering
\caption{Portfolio-level multifractal metrics}
\label{tab:portfolio}
\begin{tabular}{lc}
\toprule
\textbf{Metric} & \textbf{Value} \\
\midrule
Portfolio $\Delta h$ & 0.2405 \\
Na\"ive weighted-average $\Delta h$ & 0.2530 \\
\textbf{Multifractal Diversification Ratio (MDR)} & \textbf{1.052} \\
\bottomrule
\end{tabular}
\end{table}

The MDR is defined as:
\[
\text{MDR} = \frac{\sum_i w_i \cdot \Delta h_i}{\Delta h_{\text{portfolio}}}
\]

\begin{notebox}
MDR $> 1$ means diversification \emph{does} reduce multifractal complexity---but only by $\sim$5\%. For comparison, the classic variance-based diversification ratio for a similar portfolio would typically be 1.3--1.8$\times$. Multifractal complexity is far harder to diversify away.
\end{notebox}

% ══════════════════════════════════════════════════════════════════
\section{Individual-Asset Multifractality Impact}
% ══════════════════════════════════════════════════════════════════

\subsection{Weighted Hurst Exponents}

The weighted-average Hurst curve $\bar{h}(q) = \sum_i w_i \cdot h_i(q)$ closely tracks the actual portfolio MF-DFA curve but diverges at extreme $q$ values. This divergence is caused by \textbf{nonlinear cross-correlations} that the na\"ive weighted average ignores.

\textbf{Key observations:}
\begin{itemize}[nosep]
\item At negative $q$ (small fluctuations): BTC-USD and AAPL have the highest Hurst, pulling the portfolio toward persistence.
\item At positive $q$ (large fluctuations): JNJ and XOM dominate, making large moves more persistent than a monofractal model would predict.
\end{itemize}

\subsection{Contribution to Tail Risk (CTR)}

Using the negative-$q$ singularity spectrum ($\alpha$ for $q < 0$), we decompose which assets dominate the portfolio's left-tail fractal structure:

\[
\text{CTR}_i = \frac{w_i \cdot \bar{\alpha}_i^{(q<0)}}{\sum_j w_j \cdot \bar{\alpha}_j^{(q<0)}} \times 100\%
\]

\begin{table}[H]
\centering
\caption{Expected tail-risk contribution ranking}
\begin{tabular}{llp{7cm}}
\toprule
\textbf{Asset} & \textbf{CTR Rank} & \textbf{Rationale} \\
\midrule
BTC-USD & Highest & Widest spectrum + persistent trending ($h = 0.566$) \\
JNJ     & High    & Highest $\Delta h$ (0.308) despite low $h(2)$ \\
XOM     & High    & Second-highest $\Delta h$ (0.296) \\
EWJ     & Lowest  & Narrowest spectrum (0.185) \\
\bottomrule
\end{tabular}
\end{table}

\begin{warningbox}
JNJ and XOM are traditionally considered ``defensive'' assets. Their high CTR reveals that from a \textbf{multifractal tail-risk perspective}, they are actually among the most dangerous portfolio components. This is invisible to standard mean-variance analysis.
\end{warningbox}

\subsection{Asymmetric MF-DFA (Leverage Effect)}

By filtering MF-DFA to only upward-trending ($\text{sign\_filter} = ``+"$) or downward-trending ($\text{sign\_filter} = ``-"$) segments, we detect asymmetric persistence:

\begin{table}[H]
\centering
\caption{Asymmetric persistence patterns}
\begin{tabular}{lll}
\toprule
\textbf{Pattern} & \textbf{Meaning} & \textbf{Typical assets} \\
\midrule
$h^-(q) > h^+(q)$ & Downside more persistent & Equities (AAPL, JPM, XOM, JNJ) \\
$h^-(q) \approx h^+(q)$ & Symmetric persistence & Gold (GC=F), BTC-USD \\
$h^-(q) < h^+(q)$ & Upside more persistent (rare) & Occasionally BTC in bull regimes \\
\bottomrule
\end{tabular}
\end{table}

\begin{importantbox}
The \textbf{fractal leverage effect} means equity crashes are ``stickier'' than rallies at all timescales. Any risk measure that treats up and down moves identically will systematically underestimate drawdown risk.
\end{importantbox}

\subsection{$\Delta h \cdot w$ Heat-Map}

The heat-map of $(h_i(q) - h_{\text{portfolio}}(q)) \times w_i$ across all $q$ provides a visual fingerprint:
\begin{itemize}[nosep]
\item \textbf{Red cells} = asset increases portfolio multifractality at that $q$
\item \textbf{Blue cells} = asset decreases it
\end{itemize}

\textbf{Expected patterns:}
\begin{itemize}[nosep]
\item BTC-USD shows strong red at low $q$ (amplifies persistence of small moves)
\item GC=F and EWJ show blue (dampen multifractality---genuine diversifiers at the fractal level)
\end{itemize}

% ══════════════════════════════════════════════════════════════════
\section{Sector-Level Multifractality}
% ══════════════════════════════════════════════════════════════════

\subsection{Sector-Averaged Hurst Exponents}

\begin{table}[H]
\centering
\caption{Sector-level persistence regimes}
\begin{tabular}{lccp{5cm}}
\toprule
\textbf{Sector} & $\boldsymbol{H(2)}$ & $\boldsymbol{\Delta h}$ & \textbf{Character} \\
\midrule
Tech        & $\sim$0.50 & Moderate & Near efficient-market \\
Financials  & $\sim$0.48 & Moderate & Slightly anti-persistent \\
Energy      & $\sim$0.50 & High     & Efficient on average, complex at extremes \\
Healthcare  & $\sim$0.45 & \textbf{High} & Anti-persistent, highly multifractal \\
Commodity   & $\sim$0.43 & Moderate & Mean-reverting tendency \\
Crypto      & $\sim$\textbf{0.57} & \textbf{High} & Persistent trends, complex scaling \\
Intl ETF    & $\sim$0.43 & Low      & Closest to monofractal \\
\bottomrule
\end{tabular}
\end{table}

\subsection{Cross-Sector MF-DCCA Correlations}

The sector-pairwise $\rho_{\text{DCCA}}$ matrix measures cross-scale correlation---how synchronised two sectors' fractal structures are across all timescales.

\textbf{Expected findings:}
\begin{itemize}[nosep]
\item \textbf{High $\rho_{\text{DCCA}}$ ($>0.6$):} Tech--Financials, Energy--Industrials
\item \textbf{Low $\rho_{\text{DCCA}}$ ($<0.3$):} Gold--Equities, Crypto--Traditional
\end{itemize}

\begin{tipbox}
Sectors with high MF-DCCA correlation are likely to experience \textbf{correlated regime changes} during stress events. A crisis in Financials will propagate to Tech across all timescales---this is the fractal signature of systemic contagion.
\end{tipbox}

\subsection{Hierarchical Clustering (Dendrogram)}

The dendrogram based on MF-DCCA distances ($d = 1 - \rho_{\text{DCCA}}$) reveals a ``multifractal similarity tree.'' Expected structure:

\begin{verbatim}
  ├── Domestic equities cluster
  │   ├── Tech + Financials (tight)
  │   └── Energy + Healthcare (moderate)
  ├── International (EWJ — loosely attached)
  └── Alternative assets
      ├── Gold (GC=F)
      └── Crypto (BTC-USD)
\end{verbatim}

\begin{notebox}
This clustering is fundamentally different from correlation-based clustering. Two assets can have low Pearson correlation but high MF-DCCA if they share the same \emph{type} of multifractal dynamics. The dendrogram reveals the true ``fractal DNA'' relationships.
\end{notebox}

\subsection{Sector Tail Indices (Hill Estimator)}

The Hill tail index $\alpha$ measures how heavy the loss distribution's tail is per sector. Lower $\alpha$ = fatter tail = more extreme-event risk.

\textbf{Expected ranking} (fattest to thinnest tails):
\begin{enumerate}[nosep]
\item Crypto (lowest $\alpha$)---most extreme tail risk
\item Energy---commodity-linked fat tails
\item Tech---growth/momentum crash risk
\item Commodity (Gold)---moderate tails
\item Healthcare/Consumer---thinnest tails
\end{enumerate}

% ══════════════════════════════════════════════════════════════════
\section{Portfolio-Level Extensions}
% ══════════════════════════════════════════════════════════════════

\subsection{Volatility-Character Retention}

The $\rho_{\text{DCCA}}$ ranking between each component's $|r_i|$ and the portfolio's $|r_p|$ reveals which assets' volatility patterns survive into the aggregate:

\begin{table}[H]
\centering
\caption{Volatility-character retention ranking}
\begin{tabular}{clcc}
\toprule
\textbf{Rank} & \textbf{Asset} & \textbf{Sector} & $\boldsymbol{\overline{\rho}_{\text{DCCA}}}$ \\
\midrule
1 & JPM     & Financials & 0.782 \\
2 & EWJ     & Intl ETF   & 0.733 \\
3 & XOM     & Energy     & 0.705 \\
4 & AAPL    & Tech       & 0.689 \\
5 & JNJ     & Healthcare & 0.671 \\
6 & BTC-USD & Crypto     & 0.562 \\
7 & GC=F    & Commodity  & \textbf{0.452} \\
\bottomrule
\end{tabular}
\end{table}

\begin{importantbox}
\textbf{JPM dominates the portfolio's volatility character}---its fractal volatility pattern is almost entirely ``retained'' in the portfolio. Gold is the most genuinely diversifying asset at the fractal-volatility level, followed by BTC.
\end{importantbox}

\subsection{Rolling MDR Time-Series}

The 1-year rolling MDR tracks how multifractal diversification benefits evolve:

\begin{itemize}[nosep]
\item \textbf{MDR $> 1$} during calm markets $\to$ diversification works
\item \textbf{MDR $\to 1$ or $< 1$} during crises $\to$ diversification benefit collapses
\item Drops in MDR typically \textbf{precede or coincide} with major market events
\end{itemize}

\begin{tipbox}
Rolling MDR can serve as an \textbf{early-warning signal} for systemic risk. A sustained decline suggests that the portfolio's fractal structure is converging---assets that normally have different scaling behaviours are becoming synchronised, a hallmark of approaching crisis.
\end{tipbox}

\subsection{Multifractal-Aware Risk Budget}

Using finite-difference approximation:
\[
\text{RB}_i = w_i \cdot \frac{\partial \sigma_{\text{portfolio}}^{\text{MF-DFA}}}{\partial w_i}
\]

where $\sigma^{\text{MF-DFA}}$ is the portfolio's spectral width.

\textbf{Expected findings:}
\begin{itemize}[nosep]
\item Assets with high $\Delta h$ \textbf{and} high retention ranking (JPM, XOM) have the largest marginal risk contribution.
\item Gold (GC=F) likely has a near-zero or negative marginal contribution---increasing its weight \emph{reduces} portfolio multifractal complexity.
\item This provides \textbf{actionable weight adjustments} for portfolio construction.
\end{itemize}

\subsection{Multifractal-Adjusted Expected Shortfall}

Classical ES assumes monofractal returns ($h = 0.5$). The MF-adjusted ES scales by the empirical tail persistence:

\[
\text{ES}_{\text{adjusted}} = \text{ES}_{\text{classical}} \times \frac{h(q = -5)}{0.5}
\]

If $h(q = -5) > 0.5$ (persistent tails), ES is scaled \textbf{upward}---the classical estimate is too optimistic.

% ══════════════════════════════════════════════════════════════════
\section{Robustness Validation}
% ══════════════════════════════════════════════════════════════════

The surrogate/bootstrap hypothesis tests (199 replications) from the companion robustness notebook confirm statistical significance:

\begin{table}[H]
\centering
\caption{Surrogate hypothesis test results}
\begin{tabular}{lcccc}
\toprule
\textbf{Metric} & \textbf{Observed} & \textbf{Surrogate 95\% CI} & \textbf{$p$-value} & \textbf{Conclusion} \\
\midrule
Portfolio $|\Delta h|$ & 0.240 & $[0.009,\; 0.157]$ & \textbf{0.005} & Significant \\
AAPL $\Delta h$        & 0.232 & $[0.025,\; 0.207]$ & 0.010 & Significant \\
JPM $\Delta h$         & 0.208 & $[0.059,\; 0.227]$ & 0.135 & Borderline \\
XOM $\Delta h$         & 0.296 & $[0.025,\; 0.183]$ & \textbf{0.005} & Highly significant \\
JNJ $\Delta h$         & 0.308 & $[0.037,\; 0.209]$ & \textbf{0.005} & Highly significant \\
BTC-USD $\Delta h$     & 0.300 & $[0.057,\; 0.237]$ & \textbf{0.005} & Highly significant \\
\bottomrule
\end{tabular}
\end{table}

\begin{importantbox}
The multifractality of every asset except JPM is statistically significant beyond the 95\% confidence level. This is not a statistical artefact---the temporal structure of returns creates genuine multifractal scaling that the marginal distribution alone cannot explain.
\end{importantbox}

% ══════════════════════════════════════════════════════════════════
\section{Practical Implications}
% ══════════════════════════════════════════════════════════════════

\subsection{For Risk Management}

\begin{enumerate}[nosep]
\item \textbf{Abandon $\sqrt{T}$ scaling for VaR.} The portfolio's significant multifractality means scaling daily VaR by $\sqrt{T}$ to get monthly VaR is \textbf{guaranteed to underestimate risk} at longer horizons.

\item \textbf{Use MF-adjusted ES.} The multifractal scaling factor provides a simple correction to classical ES that accounts for empirical tail persistence.

\item \textbf{Monitor rolling MDR.} A declining MDR is an early warning that fractal diversification benefits are eroding.
\end{enumerate}

\subsection{For Portfolio Construction}

\begin{enumerate}[nosep, start=4]
\item \textbf{Fractal diversifiers matter.} Gold (GC=F) and International ETFs (EWJ) provide genuine \emph{multifractal} diversification that doesn't collapse during stress. Their low $\rho_{\text{DCCA}}$ with equities means they diversify at all timescales.

\item \textbf{Reweight using MF risk budget.} Assets with high marginal MF-DFA sensitivity (JPM, XOM) should be underweighted relative to mean-variance optimisation. Assets with low sensitivity (Gold) should be overweighted.

\item \textbf{Beware ``safe'' assets.} JNJ and XOM have the highest spectral widths despite being considered defensive. Their multifractal complexity makes them disproportionate contributors to tail risk.
\end{enumerate}

\subsection{For Research}

\begin{enumerate}[nosep, start=7]
\item \textbf{BTC-USD's persistence} ($h = 0.566$) is a genuine market inefficiency that survives surrogate testing, with implications for trend-following strategies in crypto.

\item \textbf{The fractal leverage effect} ($h^- > h^+$ for equities) suggests options pricing models should incorporate asymmetric scaling, not just asymmetric volatility.

\item \textbf{MDR $\approx$ 1.05} is a strong result---even a well-diversified 7-asset portfolio barely reduces multifractal complexity. This challenges the foundations of portfolio theory at non-Gaussian scales.
\end{enumerate}

% ══════════════════════════════════════════════════════════════════
\section{Limitations \& Next Steps}
% ══════════════════════════════════════════════════════════════════

\begin{itemize}[nosep]
\item \textbf{Small portfolio size:} 7 assets may not capture full diversification potential. The Monte Carlo analysis (200 random draws from 50 tickers) in the robustness notebook confirms MDR is stable across portfolios.

\item \textbf{Equal weighting:} Real portfolios use optimised weights. Extending to MF-DFA-constrained optimisation (maximise MDR subject to $\max(h_i(q<0)) \leq$ threshold) is a natural next step.

\item \textbf{Static analysis:} The rolling MDR captures time variation, but formal change-point detection on the rolling $h(q)$ series would provide sharper regime-transition dates.

\item \textbf{No transaction costs:} Any rebalancing strategy based on MF risk budgets should account for implementation friction.
\end{itemize}

\end{document}
